\documentclass[11pt]{article}
\usepackage{series}
\usepackage{amsmath,amssymb,amsthm}
\usepackage{geometry}
\usepackage{hyperref}

\geometry{margin=1in}

\title{\textbf{Horizons, Area Laws, and Entropy\\ from Coherence Capacity Bottlenecks}}
\author{Peter Nero}
\date{January 2026}

\begin{document}
\maketitle

\begin{abstract}
We derive horizon formation, area laws, and entropy bounds as necessary
consequences of coherence-capacity bottlenecks in projection-based
effective descriptions. Building on the theory of coherence capacity and
its transport developed in preceding work, we show that whenever inward
capacity flux exceeds outward transport, admissibility fails on a
codimension-one surface. Such surfaces act as horizons: effective
descriptions become noninvertible across them, and global information
recovery is impossible. We demonstrate that entropy arises as the
bookkeeping of coherence capacity loss and that area laws follow from
geometric bounds on capacity flux. The results are independent of
microscopic dynamics and apply to any effective theory with finite
stability margins.
\end{abstract}

\section{Introduction}

Horizons and entropy occupy a central place in modern physics.
Black holes possess entropy proportional to horizon area; quantum systems
thermalize and lose memory; effective descriptions exhibit irreversible
behavior. These phenomena are traditionally treated using separate
conceptual frameworks.

In this work we show that horizons and entropy arise from a single
structural mechanism: the formation of \emph{coherence capacity
bottlenecks}. When the transport of coherence capacity into a region
exceeds its ability to redistribute that capacity, the effective
description necessarily fails on a bounding surface. This surface acts as
a horizon for the effective description.

Entropy is then identified as the measure of coherence capacity that has
been irreversibly shed at such bottlenecks. Area laws follow from geometric
bounds on capacity flux through codimension-one surfaces.

No assumptions of fundamental information destruction, singular dynamics,
or microscopic entropy are required.

\section{Review of Coherence Capacity and Transport}

We briefly summarize the framework established previously.

\subsection{Coherence capacity}

Coherence capacity $\mathcal C(x)$ is the local stability margin that
allows a projection-based effective description to remain predictive. It
is positive inside admissible regions and vanishes at admissibility
barriers.

\subsection{Capacity transport}

Within admissible basins, coherence capacity admits a conserved current
$J^\mu_{\mathcal C}$ satisfying
\[
\nabla_\mu J^\mu_{\mathcal C} = 0 \qquad (\mathcal C > 0).
\]

Capacity is redistributed by interactions, curvature, and coupling, and
may concentrate under persistent strain.

\subsection{Capacity exhaustion}

At admissibility barriers $\mathcal C\to 0$ and capacity conservation
fails:
\[
\nabla_\mu J^\mu_{\mathcal C} = \mathcal S_B.
\]

Such barriers are associated with noninvertibility of the effective
description.
\paragraph{Scope of transport statements.}
All transport statements in this paper are formulated at the level of an effective
representation of admissibility, as introduced in the capacity transport framework. They are
slab-local and conditional on the existence of an admissible chart with controlled projection,
absolute continuity, and constitutive closure. No statement in this section asserts the
existence of a fundamental or global transport law.

\section{Capacity Bottlenecks as Horizons}

We now formalize the connection between capacity bottlenecks and horizons.

\begin{definition}[Capacity bottleneck]
A capacity bottleneck is a connected codimension-one hypersurface $H$
defined by $\mathcal C=0$, across which the effective description is
noninvertible.
\end{definition}

\begin{theorem}[Horizon formation from capacity imbalance]
If the inward flux of coherence capacity into a compact region exceeds the
maximum outward transport permitted by admissible dynamics, then a
capacity bottleneck forms on a hypersurface bounding the region.
\end{theorem}

\begin{proof}
By applying the focusing and bottleneck results established in Theorem~9.2 and
Theorem~10.1 of \emph{Dynamics of Coherence Capacity: Transport, Concentration, and Exhaustion}
to the effective current representation, persistent negative expansion along integral curves
drives $C$ to zero on a codimension-one hypersurface.

\end{proof}

\begin{corollary}
Capacity bottlenecks act as horizons: effective evolution cannot be
globally inverted across them.
\end{corollary}

This definition of a horizon is purely structural and does not rely on
light cones or causal structure. Causal horizons in geometric encodings arise as a special case of capacity bottlenecks when
the effective description admits a Lorentzian structure.


\section{Irreversibility and Loss of Global Description}

The presence of a capacity bottleneck has immediate consequences.

\begin{theorem}[No global recovery across horizons]
Across a capacity bottleneck, there exists no global inverse of the
effective evolution mapping.
\end{theorem}

\begin{proof}
This follows directly from projection noninvertibility at
$\mathcal C=0$. Distinct configurations on opposite sides of the
bottleneck share identical effective images.
\end{proof}

Thus information loss at horizons is not a failure of fundamental
dynamics, but a structural limitation of effective description.

\section{Entropy as Capacity Accounting}

We now define entropy in this framework.

\begin{definition}[Capacity entropy]
The entropy associated with a region is defined as the cumulative
coherence capacity irreversibly shed at admissibility barriers enclosing
that region.
\end{definition}

Entropy is therefore not a measure of microscopic disorder, but a measure
of lost descriptive capacity.

\begin{lemma}
Capacity entropy is nondecreasing under admissible evolution.
\end{lemma}

\begin{proof}
Capacity is conserved inside admissible regions and can only be lost at
barriers. Once lost, capacity cannot be recovered by the effective
description.
\end{proof}

This establishes a generalized second law without invoking probabilistic
assumptions.
\section{Area Laws from Capacity Flux Bounds}

We now show that capacity bottlenecks imply entropy bounds proportional to
the area of the bottleneck surface. This result is purely geometric and
does not depend on microscopic dynamics.

\subsection{Capacity flux through hypersurfaces}

Let $H$ be a capacity bottleneck hypersurface with unit normal $n_\mu$.
Define the capacity flux through $H$ by
\[
\Phi_{\mathcal C}(H) := \int_H n_\mu J^\mu_{\mathcal C}\, dA,
\]
where $dA$ is the induced area element on $H$.

Since capacity transport is local and finite-bandwidth, the flux density
$|n_\mu J^\mu_{\mathcal C}|$ admits a universal upper bound $\sigma_{\max}$
determined by admissibility constraints.

\begin{assumption}[Flux bound]
There exists a constant $\sigma_{\max}$ such that
\[
|n_\mu J^\mu_{\mathcal C}| \le \sigma_{\max}
\]
for all admissible configurations.
\end{assumption}

This assumption expresses the finite rate at which coherence capacity can
be transported across a surface.
\paragraph{Justification.}
Such a bound is the transport-level encoding of bounded projector regularity and finite
spectral separation. In admissible domains, updates of coherent structure occur at finite
effective bandwidth, which limits the rate at which admissibility can be transported across
any codimension-one interface. The flux bound therefore reflects existing MTT control data
rather than an independent physical postulate.


\subsection{Area law}

\begin{theorem}[Area bound on capacity loss]
Let $H$ be a capacity bottleneck. Then the total coherence capacity shed
across $H$ satisfies
\[
\Delta \mathcal Q_{\mathcal C}(H) \le \sigma_{\max}\, \mathrm{Area}(H).
\]
\end{theorem}

\begin{proof}
By definition,
\[
\Delta \mathcal Q_{\mathcal C}(H) = \int_H |n_\mu J^\mu_{\mathcal C}|\, dA.
\]
Applying the flux bound yields the stated inequality.
\end{proof}

\begin{corollary}[Entropy area law]
The entropy associated with a capacity bottleneck is bounded by the area
of the bottleneck surface.
\end{corollary}

Thus area laws arise as direct consequences of finite capacity flux,
independent of any holographic hypothesis.

\section{Black Hole Entropy as Capacity Saturation}

We now specialize the preceding results to black hole horizons.

\subsection{Horizons as maximal bottlenecks}

A black hole horizon corresponds to a capacity bottleneck for exterior
observables. 
\paragraph{Assumption 7.1 (Stationary saturation).}
In stationary bottleneck configurations, the effective capacity flux saturates the bound,
i.e.\ $|n_\mu J_C^\mu| = \sigma_{\max}$ on the horizon hypersurface.


\begin{theorem}[Black hole entropy]
The entropy associated with a black hole horizon is
\[
S = \sigma_{\max}\, \mathrm{Area}(H),
\]
up to universal normalization.
\end{theorem}

\begin{proof}
For stationary horizons, capacity transport is extremal. Substituting the
saturated flux into the area bound yields the stated result.
\end{proof}

This reproduces the Bekenstein--Hawking scaling without reference to
microscopic degrees of freedom or state counting.

\section{Partial Recovery and Restricted Algebras}

Although global recovery across a capacity bottleneck is impossible,
restricted recovery may remain feasible.

\subsection{Restricted observables}

Let $\mathcal A_{\mathrm{ext}}$ denote the algebra of observables supported
outside the bottleneck. Within $\mathcal A_{\mathrm{ext}}$, effective
evolution may admit a partial inverse.

\begin{theorem}[Restricted recovery]
There exists a right-inverse of the effective evolution on
$\mathcal A_{\mathrm{ext}}$ if and only if coherence capacity remains
positive for
the restriction of the coherent projection to $A_{\mathrm{ext}}$ remains bounded and admits a
measurable right inverse on the corresponding effective algebra.

\end{theorem}

\begin{proof}
Restricted observables probe only a subset of the configuration space. If
capacity is exhausted only for complementary degrees of freedom, the
projection remains invertible on the restricted sector.
\end{proof}

This explains why information recovery schemes based on restricted
observables (e.g., ``islands'') succeed without restoring global
invertibility.

\section{Relation to Holography and the Page Curve}

The coherence-capacity framework clarifies several features of modern
holographic approaches.

\subsection{Page curve}

The Page curve describes the entropy of radiation as a function of time.
In the present framework, it reflects redistribution of coherence
capacity between interior and exterior observables.

Early-time growth corresponds to capacity being shed across the horizon.
Late-time saturation reflects reorganization of capacity among restricted
algebras, not restoration of global coherence.

\subsection{No paradox}

There is no contradiction between unitary fundamental dynamics and
effective information loss. Global unitarity holds upstairs; effective
irreversibility is enforced by capacity bottlenecks.

\section{Discussion}

The results of this paper show that horizons, entropy, and area laws are
not special features of gravity but generic consequences of projection-
based effective descriptions with finite coherence capacity.

This framework explains:
\begin{itemize}
\item why entropy scales with area rather than volume,
\item why horizons imply information loss for effective observers,
\item why recovery is possible only on restricted algebras,
\item why singularities are unnecessary for horizon formation.
\end{itemize}

All results follow from finite capacity transport and do not depend on
microscopic details.

\section{Conclusion}

We have shown that coherence capacity bottlenecks give rise to horizons,
entropy, and area laws as structural necessities. Entropy measures the
coherence capacity irreversibly shed at admissibility barriers, and area
laws reflect universal bounds on capacity flux.

Together with preceding work, this completes a unified framework in which:
\begin{itemize}
\item coherence capacity is the fundamental resource,
\item gravity is its geometric bookkeeping,
\item particles and forces emerge from basin dynamics,
\item horizons and entropy arise from capacity exhaustion.
\end{itemize}

Effective physical laws persist precisely because coherence capacity is
finite; their breakdown is the price paid for stability and predictability.
\end{document}
